\documentclass{article}
% Source: Topic_05_Teacher_Edition.pdf, PDF pages 93-94 (printed pages 281A-281B): Topic Assessment Form A with the TE's marked answers, Form B with marked answers, Item Analysis for both forms. Cross-checked against no extra text source (TE images only).
\usepackage{amsmath}
\usepackage{amssymb}
\newenvironment{assessment-meta}{}{}   % topic, title, forms, item count, printed pages, sources
\newenvironment{item-analysis}{}{}     % Savvas's Item / DOK / Skills Review and Practice table (same for both forms)
\newenvironment{ta-item}[2]{}{}        % #1 = item number, #2 = form (A or B)
\newcommand{\answer}[1]{}              % the TE's marked answer, inside a ta-item
\newcommand{\placeholder}[2]{}         % #1 = type (graph, table, figure, number-line), #2 = complete description
\newcommand{\uncertain}[2]{#1}         % #1 = best reading, #2 = alternatives and why
\newcommand{\te}[1]{}                  % transcriber's note (a discrepancy, a publisher misprint, a computed check)
\begin{document}
\begin{assessment-meta}
topic: 5
title: Rational Exponents and Radical Functions
forms: A and B (parallel item for item)
items: 21 per form
printed: 281A-281B
sources: Topic_05_Teacher_Edition.pdf, PDF pages 93-94 (printed 281A Form A, 281B Form B and Item Analysis); no extra text source
\end{assessment-meta}
\te{Student page headers read "5 Topic Assessment Form A" and "5 Topic Assessment Form B"; TE running heads read "TOPIC 5 Rational Exponents and Radical Functions" and footers "TOPIC 5 | 281A | TOPIC ASSESSMENT" and "TOPIC 5 | 281B | TOPIC ASSESSMENT". They agree; no running-head misprint on these pages. Student page footers: "enVision Algebra 2 | 1 of 2 / 2 of 2 | Assessment Sourcebook".}
\te{Page 281A note box, verbatim: "Topic Assessment Assess students' understanding of topic concepts and skills using the Topic Assessment found at SavvasRealize.com. These auto-scored online assessments provide students with a breadth of technology-enhanced item types. There are two versions of the Topic Assessment, Form A and Form B. These two versions, available in print and at SavvasRealize.com, are parallel tests that assess the same content item for item. The Item Analysis chart on the next page can be used for both versions."}
% Source page 94: 281B Item Analysis
\begin{item-analysis}
Item & DOK & Skills Review and Practice \\
1 & 1 & A80 \\
2 & 2 & A80, A83 \\
3 & 2 & A80 \\
4 & 2 & F66 \\
5 & 2 & F65 \\
6 & 2 & F66 \\
7 & 2 & A80 \\
8 & 2 & F65 \\
9 & 3 & F66 \\
10 & 2 & F64 \\
11 & 2 & A81 \\
12 & 1 & F64 \\
13 & 2 & A84 \\
14 & 2 & A92 \\
15 & 2 & F65 \\
16 & 2 & A80 \\
17 & 3 & F64 \\
18 & 3 & A92 \\
19 & 3 & F64 \\
\end{item-analysis}
\te{The printed Item Analysis stops at item 19 (three columns: items 1-7, 8-14, 15-19), but both forms print 21 items. Items 20 and 21 have no row in the TE's table. The table is transcribed as printed; nothing added.}
% Source page 93: 281A Topic Assessment Form A
\begin{ta-item}{1}{A}
Simplify (\sqrt[3]{27x^9y^3}).
[ ] (x^3y)
\answer{1A. 3 (boxed: 3x^3y)}
\end{ta-item}
\begin{ta-item}{2}{A}
Multiply ((\sqrt{a+4}-2)(\sqrt{a+4}+2)).
A. (a-8)
B. (a^2)
C. (a)
D. (a^2+4a)
\answer{2A. C}
\end{ta-item}
\begin{ta-item}{3}{A}
(121^{\frac{1}{11}}) is equivalent to [box: one of]
[ ] 2
[ ] 11
[ ] (\sqrt[11]{121})
[ ] (\sqrt{121^{11}})
because [box: one of]
[ ] (11^2=121).
[ ] ((\sqrt[11]{121})^{11}=121).
[ ] (11=\sqrt{121}).
[ ] (\sqrt{121^{11}}=1).
\answer{3A. (\sqrt[11]{121}) because ((\sqrt[11]{121})^{11}=121). (third box of the first set, second box of the "because" set)}
\end{ta-item}
\begin{ta-item}{4}{A}
The graph of (y=\sqrt{x}) has been translated to the right 3 units and down 9 units. What is the equation of the translated graph?
A. (y=3+\sqrt{x+9})
B. (y=9-\sqrt{x+3})
C. (y=3-\sqrt{9-x})
D. (y=-9+\sqrt{x-3})
\answer{4A. D}
\end{ta-item}
\begin{ta-item}{5}{A}
Which expression is equivalent to ((\sqrt[3]{17})^4)?
A. (17^{\frac{4}{3}})
B. (17^{\frac{3}{4}})
C. (17^{12})
D. (\frac{4}{3}^{17})
\answer{5A. A}
\end{ta-item}
\begin{ta-item}{6}{A}
Let (f(x)=64x^3-1).
The inverse relation [box: one of]
[ ] is an inverse function.
[ ] is not an inverse function.
\answer{6A. is an inverse function.}
\end{ta-item}
\begin{ta-item}{7}{A}
Graph the function (f(x)=3+\sqrt{x}).
\placeholder{graph}{Prompt: blank coordinate grid, x-axis labeled x with arrowheads at both ends, y-axis labeled y with arrowheads at both ends; grid squares 1 unit; window x from -2 to 8 and y from -5 to 5 (grid lines); origin labeled O; x labels 2, 4, 6; y labels 4, 2, -2, -4. Printed answer (magenta, drawn curve): square-root curve starting at (0,3) on the y-axis (no endpoint dot drawn), rising and flattening to the right through about (1,4), with an arrowhead at its right end at about (4,5) where it meets the top edge of the grid.}
\answer{7A. Drawn curve: (y=3+\sqrt{x}) starting at (0,3) and rising to the right through (1,4) to (4,5), arrow at the right end.}
\end{ta-item}
\begin{ta-item}{8}{A}
Multiply (\sqrt{2}(2\sqrt{x}+\sqrt{x})).
A. (2\sqrt{3x})
B. (2\sqrt{2x}+2x)
C. (3\sqrt{2x})
D. (6\sqrt{x})
\answer{8A. C}
\end{ta-item}
\begin{ta-item}{9}{A}
Which pair of function values proves that (g(x)) is not the inverse of (f(x))?
\placeholder{table}{Two tables. Table 1: x | 2, 4, 6, 8, 10; f(x) | 6, 3, 2, 8, 4. Table 2: x | 2, 4, 6, 8, 10; g(x) | 6, 3, 2, 8, 5.}
f[ ] and g[ ]
\answer{9A. f(10) and g(4) (boxes: 10, 4)}
\te{computed: f(10)=4 but g(4)=3, not 10; the printed pair checks.}
\end{ta-item}
\begin{ta-item}{10}{A}
Which expression is equivalent to (\sqrt[3]{-x})?
A. ((-x)^3)
B. ((-x)^{\frac{1}{3}})
C. (x^{-3})
D. (x^{-\frac{1}{3}})
\answer{10A. B}
\end{ta-item}
\begin{ta-item}{11}{A}
Let (f(x)=\sqrt{x}) and (g(x)=3-x). What is the domain of (f \circ g)?
(x \le) [ ]
\answer{11A. 3 (x \le 3)}
\end{ta-item}
\begin{ta-item}{12}{A}
If (a(x)=2+\sqrt{2-8x}), what is an equation for (a^{-1}(x))?
A. (a^{-1}(x)=\frac{(x+2)^2-2}{8})
B. (a^{-1}(x)=\frac{2-(x-2)^2}{8})
C. (a^{-1}(x)=\frac{x^2-4x-8}{2})
D. (a^{-1}(x)=8x^2-4)
\answer{12A. B}
\end{ta-item}
\begin{ta-item}{13}{A}
Which expression is equivalent to (\frac{\sqrt{x}}{2\sqrt{x}-3})?
A. (\frac{5\sqrt{x}}{2x-9})
B. (\frac{5\sqrt{x}}{4x-9})
C. (\frac{2x+3\sqrt{x}}{2x-9})
D. (\frac{2x+3\sqrt{x}}{4x-9})
\answer{13A. D}
\end{ta-item}
\begin{ta-item}{14}{A}
Some values of (f(x)) are given in the table. Find the value of (f^{-1}f(10)).
\placeholder{table}{x | -6, 6, 10; f(x) | -6, 3, 6.}
(f^{-1}(f(10))=) [ ]
\answer{14A. 10}
\te{The prompt prints (f^{-1}f(10)) without the outer parentheses; the answer line prints (f^{-1}(f(10))). Both transcribed as printed.}
\end{ta-item}
\begin{ta-item}{15}{A}
A cylindrical pipe is 9 ft long and has a volume of 100 (\text{ft}^3). Find its approximate diameter to the nearest hundredth of a foot.
[ ] ft
\answer{15A. 3.76 ft}
\te{computed: (r=\sqrt{100/(9\pi)}\approx 1.881), diameter (\approx 3.76) ft; the printed key checks.}
\end{ta-item}
\begin{ta-item}{16}{A}
Solve ((x+5)^{\frac{3}{2}}=(x-1)^3).
(x=) [ ]
\answer{16A. 4}
\te{computed: (x+5=(x-1)^2) gives (x=4) or (x=-1); (x=-1) is extraneous ((4^{3/2}=8\ne -8)); the printed key checks.}
\end{ta-item}
\begin{ta-item}{17}{A}
Let (f(x)=1+\sqrt{x}). Select all the true statements.
A. (f^{-1}(-1)=0)
B. (f^{-1}(1)=0)
C. (f^{-1}(2)=1)
D. ((f \circ f^{-1})(4)=4)
E. (f^{-1}(1)\cdot f(1)=1)
\answer{17A. B, C, D}
\end{ta-item}
\begin{ta-item}{18}{A}
What property provides for the definition of (x^{\frac{a}{b}}) as (\sqrt[b]{x^a}) ?
A. Power of a Power Property
B. Product of Powers Property
C. Quotient of Powers Property
D. Square Root Property
\answer{18A. A}
\end{ta-item}
\begin{ta-item}{19}{A}
A store is having a 10\% sale on all items. There is an additional \$30 discount to the the sale price of any bicycle. Let (x) represent the price in dollars, and let (f(x)=x-30) and (g(x)=x-0.10x=0.90x) represent the discounts. Which function can the store manager use to find the final price of a bicycle?
A. (f+g)
B. (g \circ f)
C. (g-f)
D. (f \circ g)
\answer{19A. D}
\te{"to the the sale price" (doubled "the") is printed so on the student page; kept verbatim.}
\end{ta-item}
\begin{ta-item}{20}{A}
Solve (\sqrt{4x-16}=3-\sqrt{x-4}).
(x=) [ ]
\answer{20A. 5}
\end{ta-item}
\begin{ta-item}{21}{A}
The volume of a sphere is (V(r)=\frac{4}{3}\pi r^3) and the radius is increasing 2 mm per second. The function (r(t)=2t) gives the radius at time (t) seconds. Which function gives the volume at time (t)?
A. ((V \circ r)(t))
B. ((r \circ V)(t))
C. ((r+V)(t))
D. ((V \cdot r)(t))
\answer{21A. A}
\end{ta-item}
% Source page 94: 281B Topic Assessment Form B
\begin{ta-item}{1}{B}
Simplify (\sqrt{\frac{a^8}{b^{12}}}).
(\frac{a^{[ ]}}{b^{[ ]}})
\answer{1B. 4 and 6 (boxed exponents: (\frac{a^4}{b^6}))}
\end{ta-item}
\begin{ta-item}{2}{B}
Multiply ((\sqrt{x+9}-3)(\sqrt{x+9}+3)).
A. (x)
B. (x+9x)
C. (x^2)
D. (x-18)
\answer{2B. A}
\end{ta-item}
\begin{ta-item}{3}{B}
(343^{\frac{1}{7}}) is equivalent to [box: one of]
[ ] 3
[ ] (\frac{1}{49})
[ ] \uncertain{(\sqrt[\frac{1}{7}]{343})}{the root index is a tiny glyph that reads as a stacked 1 over 7 at 600 dpi; could also be a plain 1 or a dotted i; it is clearly different from the 7 index of the next choice}
[ ] (\sqrt[7]{343})
because [box: one of]
[ ] (7^3=343).
[ ] (\frac{1}{7}+\frac{2}{7}=1).
[ ] (\sqrt[7]{7^3}=\frac{1}{7^2}).
[ ] ((\sqrt[7]{343})^7=343).
\answer{3B. (\sqrt[7]{343}) because ((\sqrt[7]{343})^7=343). (fourth box of the first set, fourth box of the "because" set)}
\te{The second "because" choice prints (\frac{1}{7}+\frac{2}{7}=1), which is false as arithmetic ((3/7)); it is a distractor and is transcribed as printed.}
\end{ta-item}
\begin{ta-item}{4}{B}
The graph of (y=2+\sqrt{x}) has been translated to the left 2 units and up 1 unit. What is the equation of the translated graph?
A. (y=3+\sqrt{x+2})
B. (y=1+\sqrt{x-2})
C. (y=3-\sqrt{x+2})
D. (y=1+\sqrt{x+2})
\answer{4B. A}
\end{ta-item}
\begin{ta-item}{5}{B}
Which expression is equivalent to ((\sqrt[3]{13})^5)?
A. (13^{\frac{5}{3}})
B. (13^{\frac{3}{5}})
C. (13^{15})
D. (\frac{5}{3}^{13})
\answer{5B. A}
\end{ta-item}
\begin{ta-item}{6}{B}
Let (a(x)=-x^2+6x-10).
The inverse relation [box: one of]
[ ] is an inverse function.
[ ] is not an inverse function.
\answer{6B. is not an inverse function.}
\end{ta-item}
\begin{ta-item}{7}{B}
Graph the function (f(x)=\sqrt[3]{x-5}).
\placeholder{graph}{Prompt: blank coordinate grid, x-axis labeled x with arrowheads at both ends, y-axis labeled y with arrowheads at both ends; grid squares 1 unit; window x from -2 to 8 and y from -5 to 5 (grid lines); origin labeled O; x labels 2, 4, 6; y labels 4, 2, -2, -4. Printed answer (magenta, drawn curve): cube-root curve with arrowheads at both ends; left arrow at about (-2,-1.9) near the left edge, rising slowly through about (4,-1), steep through the inflection point (5,0) on the x-axis, then through about (6,1) and flattening to a right arrow at about (8,1.4) at the right edge.}
\answer{7B. Drawn curve: (y=\sqrt[3]{x-5}) through (4,-1), (5,0), (6,1), arrows at both ends (left end near (-2,-1.9), right end near (8,1.4)).}
\end{ta-item}
\begin{ta-item}{8}{B}
Multiply (12\sqrt{x}\cdot\sqrt{8x}).
A. (4x\sqrt{x})
B. (4\sqrt{6x})
C. (24\sqrt{2x})
D. (24x\sqrt{2})
\answer{8B. D}
\end{ta-item}
\begin{ta-item}{9}{B}
Which pair of function values proves that (g(x)) is not the inverse of (f(x))?
\placeholder{table}{Two tables. Table 1: x | 3, 6, 10, 12, 13; f(x) | 6, 3, 13, 14, 15. Table 2: x | 3, 6, 10, 12, 13; g(x) | 6, 3, 13, 11, 10.}
f[ ] and g[ ]
\answer{9B. f(13) and g(10) (boxes: 13, 10)}
\te{computed: g(10)=13 but f(13)=15, not 10; the printed pair checks.}
\end{ta-item}
\begin{ta-item}{10}{B}
Which expression is equivalent to (\sqrt[5]{-x})?
A. ((-x)^{\frac{1}{5}})
B. ((-x)^5)
C. (x^{-\frac{1}{5}})
D. (x^{-5})
\answer{10B. A}
\end{ta-item}
\begin{ta-item}{11}{B}
Let (f(x)=\sqrt{x}) and (g(x)=\frac{1}{x-5}). What is the domain of (f \circ g)?
(x \le) [ ]
\answer{11B. 5 (x \le 5)}
\te{computed: (f(g(x))=\sqrt{\frac{1}{x-5}}) is defined only for (x>5), not (x\le 5); the printed answer frame "x \le [ ]" and the key 5 are kept as printed (the frame appears to be carried over from Form A, where (x\le 3) is correct).}
\end{ta-item}
\begin{ta-item}{12}{B}
If (a(x)=1+\sqrt{2x-7}), what is an equation for (a^{-1}(x))?
A. (a^{-1}(x)=\frac{(x+1)^2+7}{14})
B. (a^{-1}(x)=\frac{(x-1)^2+7}{2})
C. (a^{-1}(x)=2(x-1)^2+7)
D. (a^{-1}(x)=x^2+2x+8)
\answer{12B. B}
\end{ta-item}
\begin{ta-item}{13}{B}
Which expression is equivalent to (\frac{2\sqrt{x}}{\sqrt{x}-3})?
A. (\frac{2+6\sqrt{x}}{x-9})
B. (\frac{2x+3\sqrt{x}}{2x-9})
C. (\frac{2x+6\sqrt{x}}{x-9})
D. (\frac{8\sqrt{x}}{2x-9})
\answer{13B. C}
\end{ta-item}
\begin{ta-item}{14}{B}
Some values of (f(x)) are given in the table. Find the value of (f^{-1}f(2)).
\placeholder{table}{x | 2, 3, 5; f(x) | 3, 1, 5.}
(f^{-1}(f(2))=) [ ]
\answer{14B. 2}
\end{ta-item}
\begin{ta-item}{15}{B}
A cylindrical pipe is 25 ft long and has a volume of 900 (\text{ft}^3). What is its approximate diameter?
A. 3.39 ft
B. 6 ft
C. 6.77 ft
D. 12 ft
\answer{15B. C}
\te{Form B item 15 is multiple choice; Form A item 15 is a fill-in box. computed: (r=\sqrt{900/(25\pi)}\approx 3.385), diameter (\approx 6.77) ft; the printed key checks (choice A is the radius).}
\end{ta-item}
\begin{ta-item}{16}{B}
Solve (x-7=\sqrt{x-5}).
(x=) [ ]
\answer{16B. -54}
\te{computed: squaring gives (x^2-15x+54=0), so (x=6) or (x=9); (x=6) is extraneous ((-1\ne 1)) and the solution is (x=9). The printed key (-54) is not a solution ((\sqrt{-59}) is undefined); it looks like the constant term of the quadratic. Printed key kept.}
\end{ta-item}
\begin{ta-item}{17}{B}
Let (f(x)=1+\sqrt{x}). Select all the true statements.
A. (f^{-1}(x)=(x-1)^2)
B. (f^{-1}(0)=1)
C. (f^{-1}(1)=4)
D. ((f^{-1}\circ f)(x)=x)
E. (f^{-1}(0)=-4)
\answer{17B. B, D}
\te{computed: (f^{-1}(x)=(x-1)^2) for (x\ge 1). B is false: (f^{-1}(0)=1) would require (f(1)=0), but (f(1)=2), and 0 is not in the range of f. A is the inverse rule (true with the domain restriction (x\ge 1), unstated). By computation only D is unconditionally true (A arguably). Printed key B, D kept.}
\end{ta-item}
\begin{ta-item}{18}{B}
What property provides for the definition of (a^{\frac{m}{n}}) as (\sqrt[n]{a^m}) ?
A. Product of Powers Property
B. Quotient of Powers Property
C. Power of a Power Property
D. Square Root Property
\answer{18B. C}
\end{ta-item}
\begin{ta-item}{19}{B}
A store increases all its prices by 20\% and then offers a \$50 discount on all purchase prices. Let (x) represent the price in dollars. Let (f(x)=1.20x) represent the increase and (g(x)=x-50) represent the discount. Which function can the store manager use to find the final purchase prices?
A. (f+g)
B. (f\times g)
C. (g \circ f)
D. (f \circ g)
\answer{19B. C}
\end{ta-item}
\begin{ta-item}{20}{B}
Solve (\sqrt{x+1}+\sqrt{x-1}=2).
(x=) [ ]
\answer{20B. 1.25}
\te{computed: (\sqrt{x-1}=\frac{1}{2}), (x=1.25); check (1.5+0.5=2); the printed key checks.}
\end{ta-item}
\begin{ta-item}{21}{B}
The volume of a sphere is (V(r)=\frac{4}{3}\pi r^3) and the radius is increasing at 5 mm per second. The function (r(t)=5t) gives the radius at time (t) seconds. Which function gives the volume at time (t)?
A. ((V \circ r)(t))
B. ((r \circ V)(t))
C. ((r+V)(t))
D. ((V \cdot r)(t))
\answer{21B. A}
\end{ta-item}
\end{document}
